%% kapa: exam=1 round=28 year=2017 period=1 form=A subject=real_estate subject_ko=부동산학원론 seq=1 number=81
%% source: Q-Net 감정평가사 기출문제 (cst00309) artlSeq=5199603 file=extracted/1008/5199603/2017년 제28회 감정평가사 1차시험 문제지/1교시 A형.pdf
%% conversion: kordoc 4.24.51 -> tools/gen1.py (auto); review=proofread-sonnet; flags=none
\begin{question}{81}
부동산정책의 공적개입 필요성에 관한 설명으로 옳지 않은 것은?
\begin{choices}
  \item 정부가 부동산시장에 개입하는 논리에는 부(-)의 외부효과 방지와 공공재 공급 등이 있다.
  \item 부동산시장은 불완전정보, 공급의 비탄력성으로 인한 수요ㆍ공급 시차로 인하여 시장실패가 나타날 수 있다.
  \item 정부는 토지를 경제적ㆍ효율적으로 이용하고 공공복리의 증진을 도모하기 위하여 용도지역제를 활용하고 있다.
  \item 정부는 주민의 편의를 위해 공공재인 도로, 공원 등의 도시계획시설을 공급하고 있다.
  \item 공공재는 시장기구에 맡겨둘 경우 경합성과 배제성으로 인하여 무임승차(free ride) 현상이 발생할 수 있다.
\end{choices}
\end{question}
\answer{5}
